\section{Introduction}
\label{sec:intro}

Feed-forward multi-view models now recover cameras and geometry from a set of images in a single forward pass~\cite{wang2024dust3r,leroy2024mast3r,wang2025vggt,wang2025pi3,lin2025da3}. VGGT~\cite{wang2025vggt}, \pithree~\cite{wang2025pi3} and Depth Anything~3 (\dathree)~\cite{lin2025da3} share one design: each view is cut into $14{\times}14$ patches, each patch becomes a token, and attention alternates between tokens of the same view and tokens of all views. The number of tokens therefore sets the cost of a run, and a practitioner with a fixed budget has to decide how to spend it: on more views, or on more pixels per view.

One answer looks especially efficient. Most of a scene can be covered by a few sharp views, with the remaining views shrunk to save tokens. Inputs of this kind also appear without anyone choosing them, when a rig combines cameras of different resolution or when photos come from different devices. Either way the model receives views of different sizes in one forward pass.

We find that the three models handle this poorly. On HiRoom, \pithree reaches 60.9 \auc when all 12 views are sharp. If 8 of those views are instead given at half resolution, accuracy falls to 0.8 (\cref{fig:teaser}). For pose, the small views can even do harm. When we score only the four sharp frames that every input shares, the mix, even with its positions aligned, is below the four sharp views alone in 28 of 30 settings, 17 of them significantly, and significantly above them in one. A mixed input is therefore not a cheaper version of a sharp one. It is a different kind of input, and the models are not built for it.

This paper asks why, and what a practitioner should do about it. We give each scene a fixed set of frames and change one factor at a time (\cref{fig:arch}). Small views can be fed as they are, with their positions mapped onto the sharp grid, with a correction to their attention, upsampled back to the sharp size, or replaced by sharp versions of the same frames. Upsampling is the key control: it adds no information, but it makes the small views look like sharp ones to the model. If upsampled views recover the loss, the loss came from how the views are represented, not from what they contain. We run the three models on five datasets, for pose and for reconstruction, at two token levels and several resolution ratios. Each study had a written protocol with predictions, fixed before any of its jobs ran. Configurations were chosen on one random seed and every reported number comes from another, a held-out dataset was never used for any choice, and the core diagnosis was repeated on a third seed with new frames. Predictions that failed are reported as failed.

Three findings come out of this.

\noindent\textbf{The failure is mostly about representation, for two of the three models.} For \pithree and \dathree, upsampling the small views recovers 62--110\% of the lost pose accuracy, although it adds nothing the small views did not already contain. \dathree gives every patch the same position in global attention, yet collapses as hard as the others, so misaligned positions across views cannot be the whole explanation. VGGT behaves differently: real detail matters to it, and so does uneven detail across views. We narrow its cause but do not establish it.

\noindent\textbf{Fixing the input works better than fixing the model.} Mapping positions onto a common grid and rescaling attention help, but recover only a small part of the loss. Resizing all views to one size, within the same token cost as the mix, recovers most of it for \pithree and \dathree.

\noindent\textbf{The appeal of mixing comes from extra views.} When tokens can be allocated freely, the best mixes do beat the best uniform inputs in some reconstruction settings. But the mixes use more views. Against uniform inputs with as many views, they win 6 and lose 11 of 54 comparisons.

\paragraph{Contributions.}
\begin{itemize}[nosep,leftmargin=*]
  \item A controlled study of mixed-resolution input for three feed-forward 3D models, five datasets, pose and reconstruction, resolution ratios from 2 to 4, three ways of splitting sharp and small views, inputs whose small views come in four different sizes, and two token levels.
  \item A diagnosis by interventions that separates lost information from representation, including a control model that rules out cross-view position misalignment as the only cause.
  \item A training-free fix that works at equal token cost, and evidence that in-model corrections do much less.
  \item An answer to the allocation question: the gain often credited to mixing resolutions is a gain from view count. We turn the results into a per-model recipe (\cref{tab:recipe}).
\end{itemize}
